% !PRINTED 484
\weberpage{484}
%<<CONT
spricht, dass die Werthe (3) von einander verschieden sind.
Permutationsgruppen $P$, die diese Eigenschaft haben, dass nämlich
die permutirten Elemente sich so in Reihen $A, B \dots S$ von
gleich vielen Elementen zerlegen lassen, dass durch keine Permu\-
tation von $P$ zwei Elemente derselben Reihe in Elemente ver\-
schiedener Reihen übergehen, heissen \gesperrt{imprimitiv}. Die einzelnen
Reihen $A, B \dots S$ heissen die \gesperrt{Systeme der Imprimitivität}.
Permutationsgruppen, bei denen eine solche Zerlegung der Ele\-
mente nicht möglich ist, heissen \gesperrt{primitive Gruppen}.

Es kann sehr wohl vorkommen, dass eine Gruppe in mehr\-
facher Art imprimitiv ist, dass sie ganz verschiedene Systeme
der Imprimitivität besitzt. So ist z. B. die Gruppe der cyklischen
Permutationen von sechs Elementen $(1, 2, 3, 4, 5, 6)$ in doppelter
Weise imprimitiv und hat die Systeme der Imprimitivität

\[
A = 1, 3, 5, \quad B = 2, 4, 6
\]
und
\[
A = 1, 4, \quad B = 2, 5, \quad C = 3, 6.
\]

Nachdem also der Begriff der primitiven und imprimitiven
Gruppen festgesetzt ist, können wir dem oben Bewiesenen den
Ausdruck geben:

1. \gesperrt{Ein imprimitiver Körper hat eine imprimitive
Gruppe.}

Es ergiebt sich aus der Imprimitivität der Gruppe für die
imprimitiven Körper ein wichtiges Resultat.

Wir wollen mit
\[
\psi(x, x_1, \dots x_{r-1})
\]
eine rationale \gesperrt{symmetrische Function} der $r$ Veränderlichen
$x, x_1 \dots x_{r-1}$ bezeichnen und setzen

\weq{4}{
\begin{aligned}
\omega &= \psi(\alpha, \alpha_1 \dots \alpha_{r-1}) \\
\omega_1 &= \psi(\beta, \beta_1 \dots \beta_{r-1}) \\
& \cdot \ \cdot \ \cdot \ \cdot \ \cdot \ \cdot \ \cdot \ \cdot \ \cdot \ \cdot \\
\omega_{s-1} &= \psi(\sigma, \sigma_1 \dots \sigma_{r-1});
\end{aligned}
}
dann ist zu beweisen, dass $\omega$ rational durch $\Theta$ ausgedrückt
werden kann, d. h. im Körper $\Omega(\Theta)$ enthalten ist. Es ist
nämlich
\weq{5}{
\varphi(t)\left(\dfrac{\omega}{t - \Theta} + \dfrac{\omega_1}{t - \Theta_1} + \cdots + \dfrac{\omega_{s-1}}{t - \Theta_{s-1}}\right) = \Phi(t)
}
%CONT>>
